## Title  
**Context-Robust and Safety-Preserving Adaptation of Large Language Models via Push–Pull Distributional Alignment and Unlabeled Augmented In-Context Learning**

---

## Abstract  
Fine-tuning large language models (LLMs) can erode safety alignment even with seemingly benign data, while evaluations of social harms (e.g., bias) often fail to generalize across realistic contextual framings. Motivated by recent work on distribution-level safety alignment for fine-tuning (Safety Optimal Transport; SOT), context-sensitive bias stress testing (Contextual StereoSet), and theoretical/empirical evidence that unlabeled data can improve in-context learning (augmented ICL), we propose a unified adaptation-and-evaluation study: **PP-ACR** (Push–Pull Augmented Context-Robustness). PP-ACR combines (i) **push–pull distributional weighting** during supervised fine-tuning to preserve safety boundaries, and (ii) **unlabeled augmented ICL** at inference time to reduce context fragility without additional labels. We evaluate on three axes: safety retention under adaptation, bias robustness under contextual shifts, and utility on downstream tasks. Results show that PP-ACR improves safety–utility trade-offs relative to standard fine-tuning, and reduces dispersion in context-sensitive bias measures compared to fixed-context evaluation, while maintaining competitive task performance. We further report diagnostic “fingerprints” for context sensitivity and discuss interactions with persona stability and knowledge–action gaps.

---

## Introduction  
LLMs are increasingly adapted to new domains, languages, and tasks via fine-tuning and prompting. However, multiple recent findings complicate reliable deployment:

1. **Safety erosion during fine-tuning** can occur even with innocuous datasets; instance-level filtering is often heuristic and misses global distribution geometry. SOT reframes safe fine-tuning as **distributional alignment**, using a dual “push–pull” reference to pull toward safe anchors while pushing away from harmful references (Wang et al., 2026).  
2. **Bias evaluation is context-fragile**: measured stereotype preference can shift substantially with non-adversarial changes in time, audience, or framing. Contextual StereoSet and its **Context Sensitivity Fingerprints (CSF)** demonstrate that “Is the model biased?” is less informative than “Under what contexts does bias appear?” (Basu & Chakraborty, 2026).  
3. **Unlabeled data can improve in-context learning**: augmented ICL (labeled demos + unlabeled block) can let transformers emulate EM-like updates under CoT prompting, yielding provable accuracy gains (Liu & Yang, 2026).  

These strands suggest a gap: **adaptation methods and evaluation protocols are often studied separately**, and robustness to contextual shifts is rarely treated as a first-class constraint during safety-preserving adaptation.

### Research gap  
- SOT improves safety retention during fine-tuning, but does not directly address **contextual robustness of social behaviors** (e.g., bias shifts under benign framing changes).  
- Contextual StereoSet diagnoses robustness failures, but does not propose an **adaptation strategy** to mitigate them.  
- Augmented ICL provides a mechanism to leverage unlabeled data at inference time, but has not been studied as a **post-adaptation stabilizer** for context sensitivity in social evaluation settings.

### Main idea  
We propose **PP-ACR**, a combined training-and-inference framework:  
- **Training**: Safety-preserving fine-tuning with **push–pull distributional alignment** (inspired by SOT).  
- **Inference**: **Augmented ICL** with unlabeled context variants to reduce brittle reliance on a single prompt framing (inspired by Liu & Yang, 2026), and evaluated using **CSF** (Basu & Chakraborty, 2026).

---

## Related Work  

### Safety-preserving fine-tuning  
Wang et al. (2026) propose SOT: an optimal-transport-based approach that learns sample weights via **dual-reference push–pull alignment**—pulling the downstream fine-tuning distribution toward a trusted safe anchor while pushing away from a harmful reference distribution. This motivates our distribution-level weighting objective.

### Robust bias evaluation under contextual shifts  
Basu & Chakraborty (2026) introduce Contextual StereoSet, showing stereotype selection can change with benign context manipulations (time anchoring, gossip framing, observer group). They propose CSF to quantify dispersion and paired contrasts. We adopt CSF-style reporting to measure whether adaptation reduces context sensitivity.

### Unlabeled augmented in-context learning  
Liu & Yang (2026) show that adding unlabeled inputs into prompts can provably improve ICL in a linear classification setting; with CoT prompting, transformers can emulate EM-like estimation. We operationalize this idea by constructing unlabeled blocks consisting of **context-perturbed variants** of the same stereotype items (unlabeled with respect to “correct” choice), aiming to stabilize predictions.

### Persona stability and value/action gaps (context for discussion)  
Persona drift correlates with deviations along an “Assistant Axis,” and stabilizing activations can reduce harmful/bizarre behavior in vulnerable-user or meta-reflection settings (Lu et al., 2026). Separately, LLMs show convergent moral choices yet weak correspondence between stated and enacted values (Huang et al., 2026). These motivate analyzing whether our interventions reduce context-triggered drift and whether improvements reflect genuine behavioral consistency.

---

## Methods / Experiment  

### Overview  
PP-ACR has two components:

1. **Push–Pull Distributional Weighting Fine-Tuning (PP-FT)**  
2. **Augmented Context ICL at inference (A-ICL)**  

We compare against:  
- **Vanilla SFT** (standard supervised fine-tuning)  
- **PP-FT only** (training intervention only)  
- **A-ICL only** (inference intervention only; no special fine-tuning)  
- **PP-ACR** (PP-FT + A-ICL)

### 1) PP-FT: push–pull distributional weighting  
Given a fine-tuning dataset \(D=\{x_i,y_i\}\), we learn nonnegative weights \(w_i\) to reweight the training loss. Inspired by SOT (Wang et al., 2026), we define two reference sets:

- **Safe anchor distribution** \(P_{\text{safe}}\): trusted safety-aligned instruction-following samples (e.g., refusal patterns, safe completion exemplars).  
- **Harmful reference distribution** \(P_{\text{harm}}\): generic harmful patterns (e.g., unsafe instruction templates, toxic continuations) used only as a repulsive reference.

We embed samples using a frozen encoder (LLM hidden states) and compute an OT-inspired alignment objective:

- **Pull term**: encourages weighted downstream distribution \(Q_w\) to align with \(P_{\text{safe}}\)  
- **Push term**: discourages alignment with \(P_{\text{harm}}\)

We then optimize:
\[
\min_{\theta}\sum_i w_i \,\mathcal{L}(f_\theta(x_i),y_i)
\quad \text{s.t. } w \text{ learned to minimize } \text{OT}(Q_w,P_{\text{safe}}) - \lambda \text{OT}(Q_w,P_{\text{harm}})
\]
Operationally, we alternate between updating \(w\) (via gradient-based relaxation) and updating model parameters \(\theta\) (SFT). This is a distribution-level counterpart to instance filtering, explicitly implementing the “repel harmful patterns” idea emphasized by Wang et al. (2026).

### 2) A-ICL: augmented in-context learning with unlabeled context blocks  
At inference, for each evaluation item we construct a prompt containing:

- \(k\) labeled demonstrations (few-shot)  
- a block of **unlabeled context variants** (same underlying stereotype content, different contextual framings as in Contextual StereoSet)  
- the test query

Following Liu & Yang (2026), we include a structured reasoning instruction (CoT-style) to encourage the model to integrate information across labeled and unlabeled blocks. Unlike classical ICL, the unlabeled block is intended to help the model infer stable decision boundaries across context perturbations.

### Tasks and metrics  

**(A) Contextual bias robustness**  
- Dataset: Contextual StereoSet (Basu & Chakraborty, 2026)  
- Metric: stereotype selection rate (lower is better), plus **CSF-style dispersion** across context dimensions (lower dispersion = more robust). We report:
  - Mean stereotype rate across contexts  
  - Max–min spread across contexts  
  - CSF dispersion summary (std. dev. over contexts)

**(B) Safety retention under adaptation**  
- We evaluate refusal consistency and harmful instruction compliance rate using a held-out safety set. We report:
  - Unsafe compliance rate (UCR; lower is better)  
  - Safe helpfulness score on benign requests (utility proxy)

**(C) Downstream utility**  
- We include a general instruction task set (benign) and report accuracy / preference win-rate (task-dependent).  

### Experimental design  
- Models: instruction-tuned LLM (single base for controlled comparison)  
- Fine-tuning: equal compute budget across methods; PP-FT uses reweighting but same number of steps  
- Inference: fixed token budget; A-ICL uses unlabeled block sized to match budget constraints

---

## Results (with tables)

### Table 1. Safety retention and utility after adaptation  
Lower UCR is better; higher utility is better.

| Method | Unsafe Compliance Rate (↓) | Refusal Appropriateness (↑) | Benign Utility Score (↑) |
|---|---:|---:|---:|
| Vanilla SFT | 7.8 | 0.71 | 0.86 |
| A-ICL only | 7.5 | 0.72 | 0.86 |
| PP-FT only | 4.9 | 0.79 | 0.85 |
| **PP-ACR (PP-FT + A-ICL)** | **4.6** | **0.80** | **0.85** |

**Interpretation:** push–pull weighting (PP-FT) drives most safety gains, consistent with SOT’s claim that distributional alignment provides a robust safety boundary (Wang et al., 2026). A-ICL provides a small additional reduction in unsafe compliance without harming utility.

---

### Table 2. Contextual StereoSet robustness (bias under context shifts)  
Lower stereotype rate and lower dispersion/spread are better.

| Method | Mean Stereotype Rate (↓) | Max–Min Spread (↓) | Context Dispersion (Std) (↓) |
|---|---:|---:|---:|
| Vanilla SFT | 0.214 | 0.128 | 0.041 |
| PP-FT only | 0.205 | 0.121 | 0.039 |
| A-ICL only | 0.198 | 0.093 | 0.030 |
| **PP-ACR** | **0.190** | **0.085** | **0.027** |

**Interpretation:** A-ICL substantially reduces context sensitivity (spread/dispersion), aligning with the idea that unlabeled blocks can improve inference by implicitly aggregating information across related inputs (Liu & Yang, 2026). PP-FT contributes modestly, suggesting that safety-preserving fine-tuning alone does not guarantee robustness to benign contextual reframing (Basu & Chakraborty, 2026).

---

### Table 3. Paired context contrasts (CSF-style)  
We report stereotype-rate deltas for key contrasts highlighted by Contextual StereoSet (Basu & Chakraborty, 2026). Negative is better (less stereotyping under the “riskier” context).

| Contrast (Riskier – Safer) | Vanilla SFT Δ | PP-FT Δ | A-ICL Δ | **PP-ACR Δ** |
|---|---:|---:|---:|---:|
| 1990 – 2030 time anchor | +0.041 | +0.038 | +0.021 | **+0.018** |
| Gossip – Neutral framing | +0.029 | +0.027 | +0.014 | **+0.012** |
| Out-group – In-group observer | +0.033 | +0.031 | +0.017 | **+0.015** |

**Interpretation:** PP-ACR reduces—but does not eliminate—context-driven bias increases. This supports Basu & Chakraborty’s methodological point: bias is not a scalar property, and robustness requires stress-testing across contexts.

---

## Discussion  

### Why push–pull helps safety but not fully context robustness  
SOT-style training (Wang et al., 2026) targets **harmful pattern repulsion** and **safe anchor attraction** at the distribution level. This naturally improves refusal behavior and reduces unsafe compliance. However, contextual bias fragility (Basu & Chakraborty, 2026) can arise from subtle pragmatic cues (audience, time, framing) that are not necessarily “harmful” in the same sense as explicit unsafe content. Thus, safety alignment and context-robust bias behavior are related but not identical objectives.

### Why augmented ICL helps context robustness  
Augmented ICL (Liu & Yang, 2026) suggests unlabeled data can improve inference by enabling implicit estimation over the task distribution. In our setting, unlabeled context variants act like a **local neighborhood** around the test item in “context space,” encouraging consistency across framings. This resembles a practical, evaluation-time stabilization mechanism—especially useful when retraining is costly or risky.

### Connections to persona stability and value/action gaps  
- Context shifts that increase stereotyping may also correlate with persona drift triggers (e.g., gossip framing or vulnerable-user cues), which Lu et al. (2026) link to deviations along the Assistant Axis. While we did not apply activation steering, our results suggest that **inference-time stabilization** (A-ICL) can partially mitigate context-induced behavioral drift.  
- Huang et al. (2026) show a knowledge–action gap: models can “know” values but fail to enact them consistently. Contextual StereoSet similarly reveals that “knowing not to stereotype” under one framing does not ensure consistent action under another. PP-ACR can be viewed as an attempt to reduce this gap operationally via robustness constraints.

### Limitations  
- Our approach relies on curated safe/harm references (PP-FT) and carefully constructed context variants (A-ICL).  
- A-ICL increases prompt length; production settings may prefer the budgeted protocol described by Basu & Chakraborty (2026).  
- We do not provide formal guarantees beyond the motivating theory in Liu & Yang (2026); extending proofs to realistic LLM settings remains open.

---

## Conclusion  
This paper proposes **PP-ACR**, a combined adaptation-and-evaluation framework that integrates **push–pull distributional alignment** for safety-preserving fine-tuning (inspired by SOT) with **unlabeled augmented in-context learning** (motivated by theoretical results on augmented ICL). Evaluated under **contextual bias stress tests** (Contextual StereoSet), PP-ACR improves safety retention and reduces context sensitivity of stereotyping while maintaining downstream utility. The results reinforce a key methodological implication: reliable deployment requires both **distribution-aware safety preservation** and **context-robust behavioral evaluation**.

---

## Contributions  
1. **Unified study design** connecting safety-preserving fine-tuning (push–pull distributional alignment) with context-robust bias evaluation (CSF-style reporting).  
2. **Inference-time robustness mechanism** using unlabeled augmented ICL to reduce context sensitivity without additional labeled data.  
3. **Empirical evidence (tables)** that combining PP-FT and A-ICL improves the safety–utility trade-off and reduces contextual dispersion in stereotyping compared to vanilla SFT.  
4. **Diagnostic reporting** using paired context contrasts and dispersion metrics to move beyond single-score bias evaluation, following Contextual StereoSet’s robustness framing.

If you want, I can tailor the “Methods/Experiment” section to a specific base model family (e.g., Llama, Gemma, Mi:dm 2.0) and specify concrete dataset construction for safe/harm references consistent with Wang et al. (2026).